% Cote 119: offprints, reports, courses, CV and letters (typed and printed part): archive cote 119, site page 162 (PDF page 163).
% typescript. Draft compared: none.
% AI-transcribed by Claude Opus 5.5 (claude-opus-5-5, adaptive thinking, effort high), Message Batches API, batch r119_full (msgbatch_01DwZNm5Y7YsKBWyAN4sFs8w), 2026-11-13; not yet checked by eye.
\scanpage{119/162}
UNIVERSITE DES SCIENCES \& TECHNIQUES \hfill Année Universitaire 1973-74

DU LANGUEDOC

\uline{Mathématiques}

M . P . II

FONCTIONS HOLOMORPHES

(Théorie de Cauchy)

0 . \uline{Introduction}

La théorie présentée dans les fascicules précédent fait apparaître
les résultats sur les fonctions analytiques comme des conséquences de la théorie
des familles sommables; Elle est donc valable dans $\mathbb{R}$ comme dans $\mathbb{C}$. Dans le
plan complexe on peut développer une théorie plus astucieuse qui donnera des
résultats souvent plus intéressants que ceux de la théorie de Weierstran. Le
résultat essentiel est qu'une fonction dérivable d'un ouvert U de $\mathbb{C}$ dans $\mathbb{C}$
est non seulement indéfiniment dérivable mais de plus analytique. (Pour la
définition de la notion de fonction dérivable se reporter au fascicule
Séries Entières. 4.5. p. 19).

1 . \uline{Prélude}. \uline{Approximation de fonctions continues par des fonctions indéfini -}
\uline{ment dérivable}.

Il est facile de construire une fonction d'un ouvert de $\mathbb{R}$ dans $\mathbb{R}$
qui soit $k$ fois dérivables mais non $k+1$ fois. On peut construire aussi une
fonction indéfiniment dérivable mais non analytique. Pour cela, d'après le
principe du prolongement analytique, il suffit de construire une fonction
indéfiniment dérivable définie sur $\mathbb{R}$ entier, non identiquement nulle et
nulle sur un ouvert de $\mathbb{R}$.

Soit
\[
\rho_{\epsilon}(x) = \begin{cases} 0 \quad \text{pour} \quad |x| \geq \epsilon \\[4pt] \dfrac{k}{\epsilon}\exp\Bigl(\dfrac{-\epsilon^{2}}{\epsilon^{2}-x^{2}}\Bigr) \quad \text{pour} \quad |x| < \epsilon \end{cases}
\]

\mspage{\hand{162}}
